\section{Conclusion}
\label{sec:conclusion}

\shb{} is a practical search mode for minimum-edge-cut compactness
optimisation.
By restarting short \recom{} chains from each burst's endpoint, it
concentrates exploration near good plans without sacrificing the
Markov property.
The implementation in \texttt{redist-cli} as \texttt{--search short-burst}
achieves 15--22\% Polsby-Popper improvement over single-seed METIS,
outperforming both multi-seed METIS and \cs{} under a 1{,}000-step
budget, and approaching \be{} performance.

The key correctness requirement is \emph{endpoint retention} rather
than within-burst minimization: retaining the burst endpoint
preserves the Markov restart property and avoids systematic bias
that would make the selected-plan sequence uninterpretable as a
well-defined stochastic process.

For the statutory baseline, \shb{} is most effective combined with
\ar{} (\texttt{--structure prime-factor}): \ar{} provides a good
initialisation, and \shb{} refines via concentrated neighbourhood
exploration.

\paragraph{Open directions.}

\begin{enumerate}
\item \textbf{Multi-run statistics.}
      This paper presents single-run results.
      Phase~2 will provide multi-run statistics (mean, standard deviation,
      and confidence intervals over 50 independent base seeds) to confirm
      that the reported improvements are consistent rather than
      seed-dependent.

\item \textbf{Combination with BisectionEnsemble.}
      The natural next step is to apply \shb{} as a post-processing
      step on the output of \be{}: \be{} optimizes each bisection node,
      and \shb{} then refines the final assembled plan.

\item \textbf{Burst length theory.}
      The empirical rule $\ell^* \approx k$ (burst length equal to
      the number of districts) warrants a formal analysis.
      The \recom{} step modifies $\approx 1/k$ of the plan per step;
      a burst of $k$ steps covers the plan once on average.
      Whether this is the optimal coverage rate for neighbourhood
      concentration is an open question.

\item \textbf{Non-uniform stationary distributions.}
      The stochastic dominance result (Proposition~\ref{thm:dominance})
      assumes a uniform stationary distribution.
      For compactness-weighted chains (where plans with higher $\pp$
      are preferred), the analysis changes: the chain naturally
      concentrates near good plans, and the Short-Burst advantage
      may be smaller or larger depending on the weight distribution.
\end{enumerate}
